
\documentclass{article}
\usepackage{colortbl}
\usepackage{makecell}
\usepackage{multirow}
\usepackage{supertabular}

\begin{document}

\newcounter{utterance}

\centering \large Interaction Transcript for game `cladder', experiment `full\_v1.5\_default', episode 8727 with reflexion{-}Qwen3.5{-}4B.
\vspace{24pt}

{ \footnotesize  \setcounter{utterance}{1}
\setlength{\tabcolsep}{0pt}
\begin{supertabular}{c@{$\;$}|p{.15\linewidth}@{}p{.15\linewidth}p{.15\linewidth}p{.15\linewidth}p{.15\linewidth}p{.15\linewidth}}
    \# & \multicolumn{2}{c}{Player} && \multicolumn{2}{c}{Game Master} \\
    \hline

    \theutterance \stepcounter{utterance}  
    & & & \multicolumn{4}{p{0.6\linewidth}}{
        \cellcolor[rgb]{0.9,0.9,0.9}{
            \makecell[{{p{\linewidth}}}]{
                \texttt{\tiny{[P1$\langle$GM]}}
                \texttt{You are an expert in causal inference. The following question is not a typical commonsense query, but rather a meticulously designed question created by a professor specializing in causal inference, intended to assess the students' mastery of the course content.} \\
\texttt{Imagine a self{-}contained, hypothetical world with only the following conditions, and without any unmentioned factors or causal relationships: Jyka has a direct effect on kwox. Yupt has a direct effect on kwox. For those who are not jyka, the probability of kwox is 32\%. For those who are jyka, the probability of kwox is 38\%. Will jyka increase the chance of kwox?} \\
\texttt{Start your answer with 'Yes' or 'No', followed by additional reasoning or evidence to support your explanation.} \\
            }
        }
    }
     \\ \\

    \theutterance \stepcounter{utterance}  
    & \multicolumn{4}{p{0.6\linewidth}}{
        \cellcolor[rgb]{0.561,0.737,0.561}{
            \makecell[{{p{\linewidth}}}]{
                \texttt{\tiny{[P1$\rangle$GM]}}
                \texttt{Yes} \\
\\ 
\texttt{**Reasoning:**} \\
\\ 
\texttt{In the context of causal inference, determining whether a variable \$X\$ (here, *jyka*) increases the chance of an outcome \$Y\$ (here, *kwox*) relies on comparing the conditional probability of the outcome given the treatment versus the control group. This is often formalized as the Average Treatment Effect (ATE) or simply by comparing \$P(Y|X=1)\$ and \$P(Y|X=0)\$.} \\
\\ 
\texttt{Based on the conditions provided in this hypothetical world:} \\
\texttt{1.  **Control Group (\$X=0\$, not *jyka*):** The probability of *kwox* is \$32\textbackslash{}\%\$.} \\
\texttt{2.  **Treatment Group (\$X=1\$, *jyka*):** The probability of *kwox* is \$38\textbackslash{}\%\$.} \\
\\ 
\texttt{To determine if *jyka* increases the chance of *kwox*, we calculate the difference:} \\
\texttt{\$\$ \textbackslash{}Delta = P(\textbackslash{}text\{kwox\} | \textbackslash{}text\{jyka\}) {-} P(\textbackslash{}text\{kwox\} | \textbackslash{}text\{not jyka\}) \$\$} \\
\texttt{\$\$ \textbackslash{}Delta = 38\textbackslash{}\% {-} 32\textbackslash{}\% = +6\textbackslash{}\% \$\$} \\
\\ 
\texttt{Since the probability of *kwox* is strictly higher in the presence of *jyka* (\$38\textbackslash{}\% > 32\textbackslash{}\%\$) and the prompt explicitly states that there are **no unmentioned factors or causal relationships** (implying no confounding variables exist to distort this observed difference), the observed association is purely causal. Therefore, the presence of *jyka* causally increases the probability of *kwox* by 6 percentage points.} \\
            }
        }
    }
    & & \\ \\

    \theutterance \stepcounter{utterance}  
    & & & \multicolumn{2}{p{0.3\linewidth}}{
        \cellcolor[rgb]{0.9,0.9,0.9}{
            \makecell[{{p{\linewidth}}}]{
                \texttt{\tiny{[GM$|$GM]}}
                \texttt{Yes} \\
            }
        }
    }
    & & \\ \\

    \theutterance \stepcounter{utterance}  
    & & & \multicolumn{2}{p{0.3\linewidth}}{
        \cellcolor[rgb]{0.9,0.9,0.9}{
            \makecell[{{p{\linewidth}}}]{
                \texttt{\tiny{[GM$|$GM]}}
                \texttt{yes} \\
            }
        }
    }
    & & \\ \\

    \theutterance \stepcounter{utterance}  
    & & & \multicolumn{2}{p{0.3\linewidth}}{
        \cellcolor[rgb]{0.9,0.9,0.9}{
            \makecell[{{p{\linewidth}}}]{
                \texttt{\tiny{[GM$|$GM]}}
                \texttt{game\_result = WIN} \\
            }
        }
    }
    & & \\ \\

\end{supertabular}
}

\end{document}
